Let \(t=\deg N=3(H_*+3W+2)\). Thus \(t=25\,308\) for uniform targets
and \(t=18\,942\) for overlapping targets. All degrees below are
multivariate total degrees. A \emph{leader} is the leading coefficient
after substituting the same indeterminate for all nine coordinates.
A nonzero such coefficient proves that the multivariate upper bound
is attained. Successor/predecessor domain shifts leave these leaders unchanged.

\begin{center}
\begin{tabular}{@{}lll@{}}
\toprule Register & Degree & Diagonal leader\\\midrule
\(N\) & \(t\) & \(2^2\)\\
\(R\) & \(3t\) & \(2^4\)\\
\(P\) & \(12t+3\) & \(2^{21}\)\\
\(A\) & \(7t+2\) & \(2^{12}\)\\
\(C\) & \(3t\) & \(2^5\)\\
\(D\) & \(20t+4\) & \(2^{34}\)\\
\(e\) & \(10t+3\) & \(2^{19}\)\\
\(L\) & \(23t+4\) & \(2^{39}\)\\
\(F\) & \(80t+18\) & \(2^{140}\)\\
\(G\) & \(103t+22\) & \(2^{179}\)\\
\(H\) & \(3t+1\) & \(2^5\)\\
\(a_0\) & \(48t+14\) & \(2^{84}\)\\
\(a_1=DFI\) & \(312t+68\) & \(2^{542}\)\\
\(\Omega\) & \(115t+3W+31\) & \(-2^{211}\)\\
\(T\), principal majorant & \(153t+33\) & \(2^{277}\)\\
\(V\), principal majorant & \(156t+34\) & \(2^{282}\)\\
\bottomrule
\end{tabular}
\end{center}

\subsection{A polynomial certificate for the hidden cancellation}
Use \(t_6,d_5,r_k,\mathcal N,L,c_*\) from the majorant proof and define
\[
 R_*=
 e(64r_1D^3+40r_2D^2+28r_3D+21r_4)+c_*D^4 .
\]
Then \(D^4T=e\,t_6\mathcal N+R_*\), and the following are exact
integer-polynomial identities:
\begin{align}
 \Delta_T:={}&D^8(T^2-\lambda^2DFG^2)\notag\\
 ={}&e^2d_5\mathcal N^2
 +1024^2D^8\bigl(X_*L^2-4DG(CD+1)\bigr)\notag\\
 &\quad+2e\,t_6\mathcal N R_*+R_*^2,\label{eq:deltaT}\\
 \Delta_V:={}&D^8(V^2-\lambda^2a_1)\notag\\
 ={}&H^2\Delta_T+\lambda^2DFD^8(H^2-1)
              +D^8(2THB_g+B_g^2).\label{eq:deltaV}
\end{align}
Here \(X_*=e^2D^3(D-1)\). The term
\(X_*\mathcal N^2-4(X_*+D)G^2\) simplifies to
\(X_*L^2-4DG(CD+1)\); this is the essential cancellation
in \eqref{eq:deltaT}.

The respective degree bounds of the four summands in
\eqref{eq:deltaT} are
\[
326t+70,\quad306t+62,\quad333t+72,\quad200t+46.
\]
For the third bound, \(\deg R_*=100t+23\), with a positive leading
coefficient supplied by \(c_*D^4\). Its term uniquely dominates.
Since \(\deg D^8=160t+32\), we have
\[
\deg(T^2-\lambda^2DFG^2)=173t+40.
\]
After multiplication by \(H^2\), that becomes \(179t+42\).
In \(V^2-\lambda^2a_1\), however, the padding cross term \(2THB_g\)
has degree \(180t+40\), which is greater because \(t>2\).
Its leader is
\(2\cdot2^{282}\cdot2^{43}=2^{326}\).
All other terms have smaller degree. Thus
\[
 \deg(V^2-\lambda^2a_1)=180t+40,\qquad
 \operatorname{lc}_{\rm diag}(V^2-\lambda^2a_1)=2^{326}.
\]

Let \(E_0=u(n+M_d)+v-2\Omega M_d\). The inner norm has the exact expansion
\[
 \Theta^2-\Omega^2(a_0+h^2b_1)
 =E_0^2-2E_0\Omega hV
       +\Omega^2h^2(V^2-b_1)-\Omega^2a_0.
\]
The third term uniquely dominates: its degree is \(410t+6W+104\)
and its diagonal leader is \(2^{748}\). Squaring produces degree
\(820t+12W+208\) and leader \(2^{1496}\).
The subtracted norm term has the \emph{same} degree and leader
\(2^{1494}\). Therefore
\[
 \boxed{\deg F_9=820t+12W+208,\qquad
        \operatorname{lc}_{\rm diag}F_9=3\cdot2^{1494}.}
\]
This proves exact degree, not merely a syntactic upper bound.
It also explains the doubled padding \(B_g\): padding only by \(B_*+1\)
would make the two highest terms cancel.

\subsection{What the certificate graph evaluates}
The primary graph evaluates exactly \(F_9\). The principal degree
certificate graph instead evaluates \(D^{16}F_9\), through
\[
 \left[
 D^8(E_0^2-2E_0\Omega hV-\Omega^2a_0)
 +\Omega^2h^2\Delta_V
 \right]^2
 -4D^{16}\Omega^4h^2a_0b_1.
\]
This avoids division while exposing every required cancellation.
Its exact degree is \(1140t+12W+272\) and its diagonal leader is
\(3\cdot2^{2038}\). Subtracting the known degree of \(D^{16}\)
recovers the primary degree. The certificate graph is not substituted
for the delivered equation.

The V8 graph uses an equal-polynomial certificate instead. Its displayed
defect has degree \(252t+56\) and diagonal leader \(2^{443}\).
Consequently
\[
 \deg F_{9,\mathrm{V8}}=964t+12W+240,\qquad
 \operatorname{lc}_{\rm diag}F_{9,\mathrm{V8}}=2^{1730}.
\]

\subsection{Unchanged-finalizer comparison and modular receipts}
For the source replacements with the inherited finalizers, the degree
formulas are
\begin{align*}
 \deg F_{9,\mathrm{inherited}}&=3456H_*+10380W+7192,\\
 \deg F_{10}&=6320H_*+19016W+12992.
\end{align*}
The inherited-finalizer graphs and their nonzero leading coefficients are
independently replayed modulo \(1\,000\,000\,007\), \(1\,000\,000\,009\), and
\(998\,244\,353\). The refined graphs use \(1\,000\,003\),
\(1\,000\,033\), and \(1\,000\,000\,007\).
A single nonzero residue at a proved multivariate upper bound suffices
for exactness; the three receipts provide additional cross-checks.
For the principal nine-variable graph the symbolic degree calculation
above also supplies an exact integer leader.

